% Generated draft: requires mathematical and source review.
\documentclass[11pt]{article}
\usepackage[margin=25mm]{geometry}
\usepackage{fontspec}
\setmainfont{Latin Modern Roman}
\usepackage{amsmath,amssymb,mathtools,unicode-math}
\setmathfont{Latin Modern Math}
\usepackage{enumitem,needspace,xurl,hyperref}
\setlength{\parindent}{0pt}
\setlength{\parskip}{5pt}
\setlength{\emergencystretch}{3em}
\newcommand{\N}{\mathbb{N}}
\newcommand{\Z}{\mathbb{Z}}
\newcommand{\Q}{\mathbb{Q}}
\newcommand{\R}{\mathbb{R}}
\newcommand{\C}{\mathbb{C}}
\newcommand{\F}{\mathbb{F}}
\newcommand{\A}{\mathbb{A}}
\newcommand{\PP}{\mathbb{P}}
\newcommand{\E}{\mathbb{E}}
\newcommand{\eps}{\varepsilon}
\newcommand{\dd}{\,\mathrm{d}}
\begin{document}
\section*{UnsolvedMath ID: 20002099}
\textbf{Source problem code:} catalogue entry
\subsection{Definitions and mathematical statement}
The item labeled Question 1.41 only records that Yi-Jen Lee discussed relating Ozsvath--Szabo and Seiberg--Witten theory. It contains no mathematical assertion or recoverable question. The aid identifies the later Kutluhan--Lee--Taubes \(HF=HM\) series as relevant context, not as the answer to an invented precise workshop conjecture.

Its additional proposition is an abstract compatibility test. Over a fixed coefficient field \(\mathbb k\), suppose two matched three-flavor exact sequences have \(\mathbb k[U]\)-module isomorphisms commuting with every exact-sequence map. Write the relevant portion as
\[
 A^\infty\xrightarrow{i_A}A^+\longrightarrow A^-,\qquad
 B^\infty\xrightarrow{i_B}B^+\longrightarrow B^-.
\]
Define the tower images and reduced quotients by
\[
 T_A=\operatorname{im}i_A,\quad A_{\mathrm{red}}=A^+/T_A,
 \qquad T_B=\operatorname{im}i_B,\quad B_{\mathrm{red}}=B^+/T_B.
\]
The report states that the commuting isomorphism carries \(T_A\) onto \(T_B\), induces an isomorphism of reduced quotients, and preserves the specified \(U\)-power filtrations, for example \(U^kA^+\) and \(U^kB^+\). In particular reduced-group vanishing is equivalent. When the usual rational-homology-sphere and coefficient hypotheses identifying such vanishing with the L-space condition apply, that condition is transported as well.

The source emphasizes limits of this formal data. A shifted tower has the same ungraded \(U\)-module but can have a different absolute correction level. A module isomorphism without distinguished-element compatibility need not identify marked contact elements. Objectwise isomorphisms of theories do not by themselves commute with cobordism maps or give a natural transformation. Absolute grading, marked classes and cobordism naturality therefore require additional explicit compatibility. These logical countermodels do not assert that the actual \(HF=HM\) constructions fail those stronger properties. The proposition is contextual synthesis rather than a new contact-geometric question or a proof of the Floer comparison itself.
\subsection{Short English statement}
The source is only a discussion note. The aid adds an abstract test showing which invariants follow from an exact-sequence-compatible \(U\)-module comparison, and which require grading, marked-element or naturality data. It does not recover a missing workshop question or limit the actual comparison theorem beyond the stated formal hypotheses.
\subsection{Sources}
\begin{itemize}
\item Local source record: catalogue entry, catalogue entry; canonical \texttt{aim-geometry-notes.json}, \texttt{notes[436]}, item 1.41. Title: ``A structured audit of what HF=HM preserves.''
\item AIM workshop list \emph{Holomorphic curves in contact geometry}, supplied reference \url{https://aimath.org/WWN/contactgeom2/contactgeom2.pdf}. Original attribution: Yi-Jen Lee's discussion of Ozsvath--Szabo and Seiberg--Witten theory. This draft uses the supplied record without accessing a PDF.
\item Supplied research summary refers to the Kutluhan--Lee--Taubes \(HF=HM\) series without full bibliographic entries. Attempt 1 supplies the compatibility proposition and shifted-tower, marked-module and objectwise-functor countermodels, whose detailed examples are not reproduced.
\item Registry dated 14 August 2026: status ``open,'' classification ``OPEN-TRIAGE,'' statement status ``exact''; attempt context only, research-run status not a problem, rigor ``proved,'' novelty confidence low. The machine-generated, non-peer-reviewed aid requires independent verification before citation.
\end{itemize}
\end{document}
